\answer{ex:09:1}{(а)~$P_\infty=0.2$, пам'ять $20\,\text{с}$. (б)~$P_\infty=1$, пам'ять $1\,\text{с}$. (в)~Пам'ять пропорційна амплітуді шуму: ушестеро.}
\answer{ex:09:2}{$P(t)=P_\infty\coth(t\sqrt{q/R})$. (а)~$\frac12\sqrt{R/q}=5\,\text{с}$ --- удвічі коротше за пам'ять. (б)~$t=\frac12\ln21\,\sqrt{R/q}\approx15\,\text{с}$: стільки має триматися підвищений \nt{ACET}.}
\answer{ex:09:3}{Дисперсія $qT/2+R/(2T)$; $T^*=\sqrt{R/q}$, дисперсія $\sqrt{qR}=P_\infty$, як у фільтра~\eqref{eq:kalman}; зростання в $(\kappa+1/\kappa)/2$ раза: $1.25$ і $5.05$.}
\answer{ex:09:4}{$a>a_w=1-\theta/(mw)$: $0.15$ при $w=0.041$, $0.22$ при $w=0.045$.}
\answer{ex:09:5}{Помилки з'являються при $M\approx40$--$60$, при $M=80$ --- $1.9$ зайвого нейрона, при $M=100$ --- частка образу $0.86$; завада від чужих образів має дисперсію $\approx(M-1)p(1-p)/N$, тож $M_{\max}\approx80$, $M_{\max}\propto N/p$.}
\answer{ex:09:6}{Наприклад, $\eta(a)=\eta\min\bigl(1,\max(0,(a-0.3)/0.6)\bigr)$. З $\eta a$ усі три чужі нейрони ввійшли в образ; з $\eta(a)$ --- жоден, запис при $a\ge0.9$ той самий (частка $1.00$).}
\answer{ex:09:7}{(а)~$40$ повторів; з $\eta(a)$ задачі~\ref{ex:09:6} --- ніколи. (б)~$200$ повторів, $10$ ночей.}
